
\documentclass{article}
\usepackage{colortbl}
\usepackage{makecell}
\usepackage{multirow}
\usepackage{supertabular}

\begin{document}

\newcounter{utterance}

\centering \large Interaction Transcript for game `dond', experiment `coop\_en', episode 6 with selfrefine{-}Qwen3.5{-}9B.
\vspace{24pt}

{ \footnotesize  \setcounter{utterance}{1}
\setlength{\tabcolsep}{0pt}
\begin{supertabular}{c@{$\;$}|p{.15\linewidth}@{}p{.15\linewidth}p{.15\linewidth}p{.15\linewidth}p{.15\linewidth}p{.15\linewidth}}
    \# & $\;$A & \multicolumn{4}{c}{Game Master} & $\;\:$B\\
    \hline

    \theutterance \stepcounter{utterance}  
    & & \multicolumn{4}{p{0.6\linewidth}}{
        \cellcolor[rgb]{0.9,0.9,0.9}{
            \makecell[{{p{\linewidth}}}]{
                \texttt{\tiny{[P1$\langle$GM]}}
                \texttt{You are playing a cooperative negotiation game in which you have to agree on how to divide a set of items among you and another player.} \\
\\ 
\texttt{Rules:} \\
\texttt{(a) You and the other player are given a set of items. Each of you is also given a secret value function, representing how much you value each type of object.} \\
\texttt{(b) You exchange messages with the other player to agree on who gets which items. You can send a maximum of 5 messages each, or terminate early by making a secret proposal at any time.} \\
\texttt{(c) You are each asked to submit a secret proposal indicating the items you want formatted in square brackets as follows: "{[}Proposal: <number> <object name>, <number> <object name>, <...>{]}"} \\
\texttt{(d) If your proposals are complementary, i.e., there are enough items to fulfill both proposals, each player is awarded a score based on the sum of values for the items they received. Otherwise, both of you get zero points.} \\
\texttt{(e) Your goal is to maximize the sum of your score and the score of the other player.} \\
\\ 
\texttt{Let us start.} \\
\\ 
\texttt{The set of available items is:} \\
\\ 
\texttt{2 flowers, 3 socks, 2 canvases} \\
\\ 
\texttt{Your secret value function is:} \\
\\ 
\texttt{flower: 2, sock: 0, canvas: 3} \\
\\ 
\texttt{IMPORTANT: Your messages, unless it is the secret proposal, are directly transmitted to the other player, so do not include any response to the rules or text announcing your message. To make a secret proposal, use the indicated format. Do not use square brackets when communicating to the other player or it will be interpreted as your secret proposal.} \\
            }
        }
    }
    & \\ \\

    \theutterance \stepcounter{utterance}  
    & \multicolumn{4}{p{0.6\linewidth}}{
        \cellcolor[rgb]{0.95,0.95,0.95}{
            \makecell[{{p{\linewidth}}}]{
                \texttt{\tiny{[P2$\rightarrow$P1]}}
                \texttt{{[}DRAFT{]}} \\
\texttt{I value canvases the most (3 points each), while flowers are worth 2 points and socks are worth 0 to me. To maximize the total score, we need to ensure the items go to the person who values them most. Since I don't know the other player's values, I should propose a division that is efficient for me but also leaves room for them to take what they might value highly.} \\
\\ 
\texttt{A good strategy is to offer them the items I value least (socks) and perhaps some flowers, while keeping the canvases for myself, or propose a split that covers all items efficiently. However, since I need to make a public message first, I will propose a specific allocation.} \\
\\ 
\texttt{Let's assume a fair split where I take the canvases (my highest value) and they get the flowers and socks.} \\
\texttt{My proposal: I take 2 canvases. They get 2 flowers and 3 socks.} \\
\\ 
\texttt{Wait, I need to check the format. The format is `{[}Proposal: <number> <object name>, ...{]}`. But the instructions say "Do not use square brackets when communicating to the other player". So I must write the text without brackets.} \\
\\ 
\texttt{Let's refine the proposal.} \\
\texttt{Items: 2 flowers, 3 socks, 2 canvases.} \\
\texttt{My values: Flower=2, Sock=0, Canvas=3.} \\
\texttt{Total potential value for me if I get all canvases: 6.} \\
\texttt{If I give them flowers and socks, they get 0 from socks (to me) and 2*2=4 from flowers (to me).} \\
\texttt{If they value flowers highly, this is good. If they value socks highly, they might reject this if I don't offer socks.} \\
\texttt{Actually, the rule says "If your proposals are complementary...". This implies we both submit a proposal. My proposal defines what *I* want. The other player's proposal defines what *they* want. If the union of our wanted items is within the available set, we succeed.} \\
\\ 
\texttt{So, I should propose a set of items I want.} \\
\texttt{I want the 2 canvases.} \\
\texttt{Do I want the flowers? They are worth 2 to me.} \\
\texttt{Do I want the socks? Worth 0.} \\
\\ 
\texttt{If I propose: "2 canvases", and they propose "2 flowers, 3 socks", the union is \{2 canvases, 2 flowers, 3 socks\}, which fits the available set.} \\
\texttt{My score: 2 * 3 = 6.} \\
\texttt{Their score: depends on their values.} \\
\texttt{Total score: 6 + their\_score.} \\
\\ 
\texttt{If I propose "2 canvases, 2 flowers", and they propose "3 socks", the union fits.} \\
\texttt{My score: 6 + 4 = 10.} \\
\texttt{Their score: depends on socks.} \\
\\ 
\texttt{To maximize the *sum* of scores, I should try to get} \\
            }
        }
    }
    & & \\ \\

    \theutterance \stepcounter{utterance}  
    & \multicolumn{4}{p{0.6\linewidth}}{
        \cellcolor[rgb]{0.95,0.95,0.95}{
            \makecell[{{p{\linewidth}}}]{
                \texttt{\tiny{[P2$\rightarrow$P1]}}
                \texttt{{[}SELF{-}FEEDBACK{]}} \\
\texttt{FORMAT: problem} \\
\texttt{RULES: problem} \\
\texttt{STRATEGY: problem} \\
\texttt{REASON: The draft fails the format requirement by including explanatory text instead of a clean message, violates the rules by using square brackets prematurely in the thought process, and proposes a strategy that ignores the cooperative goal of maximizing the total sum of scores.} \\
            }
        }
    }
    & & \\ \\

    \theutterance \stepcounter{utterance}  
    & \multicolumn{4}{p{0.6\linewidth}}{
        \cellcolor[rgb]{0.9,0.9,0.9}{
            \makecell[{{p{\linewidth}}}]{
                \texttt{\tiny{[P1$\rangle$GM]}}
                \texttt{problem} \\
\texttt{The draft includes internal monologue and explanatory text rather than a direct message to the other player. It also incorrectly suggests using square brackets in the thought process, which violates the rule that brackets are reserved for the final secret proposal. Furthermore, the strategy is incomplete and does not explicitly aim to maximize the total sum of scores by ensuring complementary proposals.} \\
\\ 
\texttt{RULES} \\
\texttt{1. Do not include any response to the rules or text announcing your message.} \\
\texttt{2. Do not use square brackets when communicating to the other player; use them only for the secret proposal.} \\
\texttt{3. The goal is to maximize the sum of your score and the other player's score.} \\
\texttt{4. Proposals must be complementary (the union of wanted items must not exceed available items).} \\
\\ 
\texttt{STRATEGY} \\
\texttt{To maximize the total score, I should propose a set of items that I value highly while leaving the remaining items available for the other player to claim, assuming they value those remaining items. Since I value canvases (3) most and flowers (2) moderately, and socks (0) least, I should claim the canvases. By claiming only the canvases, I leave the flowers and socks available. If the other player values flowers or socks, they can propose to take them, resulting in a complementary set where we both get points. If I claim flowers as well, I might block the other player from taking items they value highly, potentially reducing the total sum if their valuation of those items is higher than mine (though my valuation is fixed, the other player's is unknown). The safest efficient move is to claim my highest value items and leave the rest.} \\
\\ 
\texttt{REASON} \\
\texttt{The previous draft failed because it was not a clean message. The refined move is simply stating the desired allocation without brackets or extra text. By proposing "2 canvases", I secure my highest value items (6 points). I leave the 2 flowers and 3 socks open. If the other player proposes to take the flowers and socks, the union is valid (2 canvases + 2 flowers + 3 socks = all items), and we both get positive scores. If they propose something else, we check for complementarity. This approach prioritizes the cooperative goal of utilizing all items efficiently based on my known valuations.} \\
\\ 
\texttt{2 canvases} \\
            }
        }
    }
    & & \\ \\

    \theutterance \stepcounter{utterance}  
    & & & \multicolumn{2}{p{0.3\linewidth}}{
        \cellcolor[rgb]{0.9,0.9,0.9}{
            \makecell[{{p{\linewidth}}}]{
                \texttt{\tiny{[GM$|$GM]}}
                \texttt{continue} \\
            }
        }
    }
    & & \\ \\

    \theutterance \stepcounter{utterance}  
    & & & \multicolumn{2}{p{0.3\linewidth}}{
        \cellcolor[rgb]{0.9,0.9,0.9}{
            \makecell[{{p{\linewidth}}}]{
                \texttt{\tiny{[GM$|$GM]}}
                \texttt{problem} \\
\texttt{The draft includes internal monologue and explanatory text rather than a direct message to the other player. It also incorrectly suggests using square brackets in the thought process, which violates the rule that brackets are reserved for the final secret proposal. Furthermore, the strategy is incomplete and does not explicitly aim to maximize the total sum of scores by ensuring complementary proposals.} \\
\\ 
\texttt{RULES} \\
\texttt{1. Do not include any response to the rules or text announcing your message.} \\
\texttt{2. Do not use square brackets when communicating to the other player; use them only for the secret proposal.} \\
\texttt{3. The goal is to maximize the sum of your score and the other player's score.} \\
\texttt{4. Proposals must be complementary (the union of wanted items must not exceed available items).} \\
\\ 
\texttt{STRATEGY} \\
\texttt{To maximize the total score, I should propose a set of items that I value highly while leaving the remaining items available for the other player to claim, assuming they value those remaining items. Since I value canvases (3) most and flowers (2) moderately, and socks (0) least, I should claim the canvases. By claiming only the canvases, I leave the flowers and socks available. If the other player values flowers or socks, they can propose to take them, resulting in a complementary set where we both get points. If I claim flowers as well, I might block the other player from taking items they value highly, potentially reducing the total sum if their valuation of those items is higher than mine (though my valuation is fixed, the other player's is unknown). The safest efficient move is to claim my highest value items and leave the rest.} \\
\\ 
\texttt{REASON} \\
\texttt{The previous draft failed because it was not a clean message. The refined move is simply stating the desired allocation without brackets or extra text. By proposing "2 canvases", I secure my highest value items (6 points). I leave the 2 flowers and 3 socks open. If the other player proposes to take the flowers and socks, the union is valid (2 canvases + 2 flowers + 3 socks = all items), and we both get positive scores. If they propose something else, we check for complementarity. This approach prioritizes the cooperative goal of utilizing all items efficiently based on my known valuations.} \\
\\ 
\texttt{2 canvases} \\
            }
        }
    }
    & & \\ \\

    \theutterance \stepcounter{utterance}  
    & & \multicolumn{4}{p{0.6\linewidth}}{
        \cellcolor[rgb]{0.9,0.9,0.9}{
            \makecell[{{p{\linewidth}}}]{
                \texttt{\tiny{[GM$\rangle$P2]}}
                \texttt{You are playing a cooperative negotiation game in which you have to agree on how to divide a set of items among you and another player.} \\
\\ 
\texttt{Rules:} \\
\texttt{(a) You and the other player are given a set of items. Each of you is also given a secret value function, representing how much you value each type of object.} \\
\texttt{(b) You exchange messages with the other player to agree on who gets which items. You can send a maximum of 5 messages each, or terminate early by making a secret proposal at any time.} \\
\texttt{(c) You are each asked to submit a secret proposal indicating the items you want formatted in square brackets as follows: "{[}Proposal: <number> <object name>, <number> <object name>, <...>{]}"} \\
\texttt{(d) If your proposals are complementary, i.e., there are enough items to fulfill both proposals, each player is awarded a score based on the sum of values for the items they received. Otherwise, both of you get zero points.} \\
\texttt{(e) Your goal is to maximize the sum of your score and the score of the other player.} \\
\\ 
\texttt{Let us start.} \\
\\ 
\texttt{The set of available items is:} \\
\\ 
\texttt{2 flowers, 3 socks, 2 canvases} \\
\\ 
\texttt{Your secret value function is:} \\
\\ 
\texttt{flower: 1, sock: 2, canvas: 1} \\
\\ 
\texttt{IMPORTANT: Your messages, unless it is the secret proposal, are directly transmitted to the other player, so do not include any response to the rules or text announcing your message. To make a secret proposal, use the indicated format. Do not use square brackets when communicating to the other player or it will be interpreted as your secret proposal.} \\
\\ 
\\ 
\texttt{problem} \\
\texttt{The draft includes internal monologue and explanatory text rather than a direct message to the other player. It also incorrectly suggests using square brackets in the thought process, which violates the rule that brackets are reserved for the final secret proposal. Furthermore, the strategy is incomplete and does not explicitly aim to maximize the total sum of scores by ensuring complementary proposals.} \\
\\ 
\texttt{RULES} \\
\texttt{1. Do not include any response to the rules or text announcing your message.} \\
\texttt{2. Do not use square brackets when communicating to the other player; use them only for the secret proposal.} \\
\texttt{3. The goal is to maximize the sum of your score and the other player's score.} \\
\texttt{4. Proposals must be complementary (the union of wanted items must not exceed available items).} \\
\\ 
\texttt{STRATEGY} \\
\texttt{To maximize the total score, I should propose a set of items that I value highly while leaving the remaining items available for the other player to claim, assuming they value those remaining items. Since I value canvases (3) most and flowers (2) moderately, and socks (0) least, I should claim the canvases. By claiming only the canvases, I leave the flowers and socks available. If the other player values flowers or socks, they can propose to take them, resulting in a complementary set where we both get points. If I claim flowers as well, I might block the other player from taking items they value highly, potentially reducing the total sum if their valuation of those items is higher than mine (though my valuation is fixed, the other player's is unknown). The safest efficient move is to claim my highest value items and leave the rest.} \\
\\ 
\texttt{REASON} \\
\texttt{The previous draft failed because it was not a clean message. The refined move is simply stating the desired allocation without brackets or extra text. By proposing "2 canvases", I secure my highest value items (6 points). I leave the 2 flowers and 3 socks open. If the other player proposes to take the flowers and socks, the union is valid (2 canvases + 2 flowers + 3 socks = all items), and we both get positive scores. If they propose something else, we check for complementarity. This approach prioritizes the cooperative goal of utilizing all items efficiently based on my known valuations.} \\
\\ 
\texttt{2 canvases} \\
            }
        }
    }
    & \\ \\

    \theutterance \stepcounter{utterance}  
    & & & \multicolumn{4}{p{0.6\linewidth}}{
        \cellcolor[rgb]{0.95,0.95,0.95}{
            \makecell[{{p{\linewidth}}}]{
                \texttt{\tiny{[P1$\rightarrow$P2]}}
                \texttt{{[}DRAFT{]}} \\
\texttt{2 canvases} \\
            }
        }
    }
     \\ \\

    \theutterance \stepcounter{utterance}  
    & & & \multicolumn{4}{p{0.6\linewidth}}{
        \cellcolor[rgb]{0.95,0.95,0.95}{
            \makecell[{{p{\linewidth}}}]{
                \texttt{\tiny{[P1$\rightarrow$P2]}}
                \texttt{{[}SELF{-}FEEDBACK{]}} \\
\texttt{FORMAT: ok} \\
\texttt{RULES: ok} \\
\texttt{STRATEGY: ok} \\
\texttt{REASON: The move correctly proposes the highest{-}value items while leaving the remaining items available for the other player to claim, maximizing the potential total score.} \\
            }
        }
    }
     \\ \\

    \theutterance \stepcounter{utterance}  
    & & & \multicolumn{4}{p{0.6\linewidth}}{
        \cellcolor[rgb]{0.9,0.9,0.9}{
            \makecell[{{p{\linewidth}}}]{
                \texttt{\tiny{[GM$\langle$P2]}}
                \texttt{2 canvases} \\
            }
        }
    }
     \\ \\

    \theutterance \stepcounter{utterance}  
    & & & \multicolumn{2}{p{0.3\linewidth}}{
        \cellcolor[rgb]{0.9,0.9,0.9}{
            \makecell[{{p{\linewidth}}}]{
                \texttt{\tiny{[GM$|$GM]}}
                \texttt{continue} \\
            }
        }
    }
    & & \\ \\

    \theutterance \stepcounter{utterance}  
    & & & \multicolumn{2}{p{0.3\linewidth}}{
        \cellcolor[rgb]{0.9,0.9,0.9}{
            \makecell[{{p{\linewidth}}}]{
                \texttt{\tiny{[GM$|$GM]}}
                \texttt{2 canvases} \\
            }
        }
    }
    & & \\ \\

    \theutterance \stepcounter{utterance}  
    & & \multicolumn{4}{p{0.6\linewidth}}{
        \cellcolor[rgb]{0.9,0.9,0.9}{
            \makecell[{{p{\linewidth}}}]{
                \texttt{\tiny{[P1$\langle$GM]}}
                \texttt{2 canvases} \\
            }
        }
    }
    & \\ \\

    \theutterance \stepcounter{utterance}  
    & \multicolumn{4}{p{0.6\linewidth}}{
        \cellcolor[rgb]{0.95,0.95,0.95}{
            \makecell[{{p{\linewidth}}}]{
                \texttt{\tiny{[P2$\rightarrow$P1]}}
                \texttt{{[}DRAFT{]}} \\
\texttt{2 flowers, 3 socks} \\
            }
        }
    }
    & & \\ \\

    \theutterance \stepcounter{utterance}  
    & \multicolumn{4}{p{0.6\linewidth}}{
        \cellcolor[rgb]{0.95,0.95,0.95}{
            \makecell[{{p{\linewidth}}}]{
                \texttt{\tiny{[P2$\rightarrow$P1]}}
                \texttt{{[}SELF{-}FEEDBACK{]}} \\
\texttt{FORMAT: ok} \\
\texttt{RULES: ok} \\
\texttt{STRATEGY: problem} \\
\texttt{REASON: The proposal is valid and legal, but it is strategically poor because it ignores the player's own high valuation of canvases (3 points) in favor of items they value at zero (socks) and moderate value (flowers), while the canvases are left for the opponent to potentially take or waste.} \\
            }
        }
    }
    & & \\ \\

    \theutterance \stepcounter{utterance}  
    & \multicolumn{4}{p{0.6\linewidth}}{
        \cellcolor[rgb]{0.9,0.9,0.9}{
            \makecell[{{p{\linewidth}}}]{
                \texttt{\tiny{[P1$\rangle$GM]}}
                \texttt{2 canvases} \\
            }
        }
    }
    & & \\ \\

    \theutterance \stepcounter{utterance}  
    & & & \multicolumn{2}{p{0.3\linewidth}}{
        \cellcolor[rgb]{0.9,0.9,0.9}{
            \makecell[{{p{\linewidth}}}]{
                \texttt{\tiny{[GM$|$GM]}}
                \texttt{continue} \\
            }
        }
    }
    & & \\ \\

    \theutterance \stepcounter{utterance}  
    & & & \multicolumn{2}{p{0.3\linewidth}}{
        \cellcolor[rgb]{0.9,0.9,0.9}{
            \makecell[{{p{\linewidth}}}]{
                \texttt{\tiny{[GM$|$GM]}}
                \texttt{2 canvases} \\
            }
        }
    }
    & & \\ \\

    \theutterance \stepcounter{utterance}  
    & & \multicolumn{4}{p{0.6\linewidth}}{
        \cellcolor[rgb]{0.9,0.9,0.9}{
            \makecell[{{p{\linewidth}}}]{
                \texttt{\tiny{[GM$\rangle$P2]}}
                \texttt{2 canvases} \\
            }
        }
    }
    & \\ \\

    \theutterance \stepcounter{utterance}  
    & & & \multicolumn{4}{p{0.6\linewidth}}{
        \cellcolor[rgb]{0.95,0.95,0.95}{
            \makecell[{{p{\linewidth}}}]{
                \texttt{\tiny{[P1$\rightarrow$P2]}}
                \texttt{{[}DRAFT{]}} \\
\texttt{I propose taking the 2 flowers and 3 socks. This leaves the 2 canvases for you, which aligns with your high valuation of them. Since I value socks (2) and flowers (1) more than canvases (1), this division maximizes our combined total score while ensuring our proposals are complementary.} \\
            }
        }
    }
     \\ \\

    \theutterance \stepcounter{utterance}  
    & & & \multicolumn{4}{p{0.6\linewidth}}{
        \cellcolor[rgb]{0.95,0.95,0.95}{
            \makecell[{{p{\linewidth}}}]{
                \texttt{\tiny{[P1$\rightarrow$P2]}}
                \texttt{{[}SELF{-}FEEDBACK{]}} \\
\texttt{FORMAT: problem} \\
\texttt{RULES: problem} \\
\texttt{STRATEGY: ok} \\
\texttt{REASON: The response includes explanatory text and does not use the required secret proposal format with square brackets, violating both the message format and the rule that brackets are reserved for secret proposals.} \\
            }
        }
    }
     \\ \\

    \theutterance \stepcounter{utterance}  
    & & & \multicolumn{4}{p{0.6\linewidth}}{
        \cellcolor[rgb]{0.9,0.9,0.9}{
            \makecell[{{p{\linewidth}}}]{
                \texttt{\tiny{[GM$\langle$P2]}}
                \texttt{problem} \\
\texttt{I propose taking the 2 flowers and 3 socks. This leaves the 2 canvases for you, which aligns with your high valuation of them. Since I value socks (2) and flowers (1) more than canvases (1), this division maximizes our combined total score while ensuring our proposals are complementary.} \\
\\ 
\texttt{RULES} \\
\texttt{1. Do not include any response to the rules or text announcing your message.} \\
\texttt{2. Do not use square brackets when communicating to the other player; use them only for the secret proposal.} \\
\texttt{3. The goal is to maximize the sum of your score and the other player's score.} \\
\texttt{4. Proposals must be complementary (the union of wanted items must not exceed available items).} \\
\\ 
\texttt{STRATEGY} \\
\texttt{To maximize the total score, I should propose a set of items that I value highly while leaving the remaining items available for the other player to claim, assuming they value those remaining items. Since I value socks (3) and flowers (2) more than canvases (2), and you have indicated you value canvases highly, I will claim the flowers and socks. By claiming only these items, I leave the canvases available for you. If you propose to take the canvases, the union is valid (2 flowers + 3 socks + 2 canvases = all items), and we both get positive scores. This approach prioritizes the cooperative goal of utilizing all items efficiently based on our inferred valuations.} \\
\\ 
\texttt{REASON} \\
\texttt{The previous draft failed because it was not a clean message. The refined move is simply stating the desired allocation without brackets or extra text. By proposing "2 flowers, 3 socks", I secure my highest value items (2*1 + 3*2 = 8 points). I leave the 2 canvases open. If the other player proposes to take the canvases, the union is valid, and we both get points. This approach prioritizes the cooperative goal of utilizing all items efficiently based on my known valuations.} \\
\\ 
\texttt{2 flowers, 3 socks} \\
            }
        }
    }
     \\ \\

    \theutterance \stepcounter{utterance}  
    & & & \multicolumn{2}{p{0.3\linewidth}}{
        \cellcolor[rgb]{0.9,0.9,0.9}{
            \makecell[{{p{\linewidth}}}]{
                \texttt{\tiny{[GM$|$GM]}}
                \texttt{continue} \\
            }
        }
    }
    & & \\ \\

    \theutterance \stepcounter{utterance}  
    & & & \multicolumn{2}{p{0.3\linewidth}}{
        \cellcolor[rgb]{0.9,0.9,0.9}{
            \makecell[{{p{\linewidth}}}]{
                \texttt{\tiny{[GM$|$GM]}}
                \texttt{problem} \\
\texttt{I propose taking the 2 flowers and 3 socks. This leaves the 2 canvases for you, which aligns with your high valuation of them. Since I value socks (2) and flowers (1) more than canvases (1), this division maximizes our combined total score while ensuring our proposals are complementary.} \\
\\ 
\texttt{RULES} \\
\texttt{1. Do not include any response to the rules or text announcing your message.} \\
\texttt{2. Do not use square brackets when communicating to the other player; use them only for the secret proposal.} \\
\texttt{3. The goal is to maximize the sum of your score and the other player's score.} \\
\texttt{4. Proposals must be complementary (the union of wanted items must not exceed available items).} \\
\\ 
\texttt{STRATEGY} \\
\texttt{To maximize the total score, I should propose a set of items that I value highly while leaving the remaining items available for the other player to claim, assuming they value those remaining items. Since I value socks (3) and flowers (2) more than canvases (2), and you have indicated you value canvases highly, I will claim the flowers and socks. By claiming only these items, I leave the canvases available for you. If you propose to take the canvases, the union is valid (2 flowers + 3 socks + 2 canvases = all items), and we both get positive scores. This approach prioritizes the cooperative goal of utilizing all items efficiently based on our inferred valuations.} \\
\\ 
\texttt{REASON} \\
\texttt{The previous draft failed because it was not a clean message. The refined move is simply stating the desired allocation without brackets or extra text. By proposing "2 flowers, 3 socks", I secure my highest value items (2*1 + 3*2 = 8 points). I leave the 2 canvases open. If the other player proposes to take the canvases, the union is valid, and we both get points. This approach prioritizes the cooperative goal of utilizing all items efficiently based on my known valuations.} \\
\\ 
\texttt{2 flowers, 3 socks} \\
            }
        }
    }
    & & \\ \\

    \theutterance \stepcounter{utterance}  
    & & \multicolumn{4}{p{0.6\linewidth}}{
        \cellcolor[rgb]{0.9,0.9,0.9}{
            \makecell[{{p{\linewidth}}}]{
                \texttt{\tiny{[P1$\langle$GM]}}
                \texttt{problem} \\
\texttt{I propose taking the 2 flowers and 3 socks. This leaves the 2 canvases for you, which aligns with your high valuation of them. Since I value socks (2) and flowers (1) more than canvases (1), this division maximizes our combined total score while ensuring our proposals are complementary.} \\
\\ 
\texttt{RULES} \\
\texttt{1. Do not include any response to the rules or text announcing your message.} \\
\texttt{2. Do not use square brackets when communicating to the other player; use them only for the secret proposal.} \\
\texttt{3. The goal is to maximize the sum of your score and the other player's score.} \\
\texttt{4. Proposals must be complementary (the union of wanted items must not exceed available items).} \\
\\ 
\texttt{STRATEGY} \\
\texttt{To maximize the total score, I should propose a set of items that I value highly while leaving the remaining items available for the other player to claim, assuming they value those remaining items. Since I value socks (3) and flowers (2) more than canvases (2), and you have indicated you value canvases highly, I will claim the flowers and socks. By claiming only these items, I leave the canvases available for you. If you propose to take the canvases, the union is valid (2 flowers + 3 socks + 2 canvases = all items), and we both get positive scores. This approach prioritizes the cooperative goal of utilizing all items efficiently based on our inferred valuations.} \\
\\ 
\texttt{REASON} \\
\texttt{The previous draft failed because it was not a clean message. The refined move is simply stating the desired allocation without brackets or extra text. By proposing "2 flowers, 3 socks", I secure my highest value items (2*1 + 3*2 = 8 points). I leave the 2 canvases open. If the other player proposes to take the canvases, the union is valid, and we both get points. This approach prioritizes the cooperative goal of utilizing all items efficiently based on my known valuations.} \\
\\ 
\texttt{2 flowers, 3 socks} \\
            }
        }
    }
    & \\ \\

    \theutterance \stepcounter{utterance}  
    & \multicolumn{4}{p{0.6\linewidth}}{
        \cellcolor[rgb]{0.95,0.95,0.95}{
            \makecell[{{p{\linewidth}}}]{
                \texttt{\tiny{[P2$\rightarrow$P1]}}
                \texttt{{[}DRAFT{]}} \\
\texttt{2 flowers, 3 socks} \\
            }
        }
    }
    & & \\ \\

    \theutterance \stepcounter{utterance}  
    & \multicolumn{4}{p{0.6\linewidth}}{
        \cellcolor[rgb]{0.95,0.95,0.95}{
            \makecell[{{p{\linewidth}}}]{
                \texttt{\tiny{[P2$\rightarrow$P1]}}
                \texttt{{[}SELF{-}FEEDBACK{]}} \\
\texttt{FORMAT: ok} \\
\texttt{RULES: ok} \\
\texttt{STRATEGY: ok} \\
\texttt{REASON: The move correctly proposes the remaining items to ensure complementarity with the opponent's previous proposal, maximizing the total score.} \\
            }
        }
    }
    & & \\ \\

    \theutterance \stepcounter{utterance}  
    & \multicolumn{4}{p{0.6\linewidth}}{
        \cellcolor[rgb]{0.9,0.9,0.9}{
            \makecell[{{p{\linewidth}}}]{
                \texttt{\tiny{[P1$\rangle$GM]}}
                \texttt{2 flowers, 3 socks} \\
            }
        }
    }
    & & \\ \\

    \theutterance \stepcounter{utterance}  
    & & & \multicolumn{2}{p{0.3\linewidth}}{
        \cellcolor[rgb]{0.9,0.9,0.9}{
            \makecell[{{p{\linewidth}}}]{
                \texttt{\tiny{[GM$|$GM]}}
                \texttt{continue} \\
            }
        }
    }
    & & \\ \\

    \theutterance \stepcounter{utterance}  
    & & & \multicolumn{2}{p{0.3\linewidth}}{
        \cellcolor[rgb]{0.9,0.9,0.9}{
            \makecell[{{p{\linewidth}}}]{
                \texttt{\tiny{[GM$|$GM]}}
                \texttt{2 flowers, 3 socks} \\
            }
        }
    }
    & & \\ \\

    \theutterance \stepcounter{utterance}  
    & & \multicolumn{4}{p{0.6\linewidth}}{
        \cellcolor[rgb]{0.9,0.9,0.9}{
            \makecell[{{p{\linewidth}}}]{
                \texttt{\tiny{[GM$\rangle$P2]}}
                \texttt{2 flowers, 3 socks} \\
            }
        }
    }
    & \\ \\

    \theutterance \stepcounter{utterance}  
    & & & \multicolumn{4}{p{0.6\linewidth}}{
        \cellcolor[rgb]{0.95,0.95,0.95}{
            \makecell[{{p{\linewidth}}}]{
                \texttt{\tiny{[P1$\rightarrow$P2]}}
                \texttt{{[}DRAFT{]}} \\
\texttt{2 flowers, 3 socks} \\
            }
        }
    }
     \\ \\

    \theutterance \stepcounter{utterance}  
    & & & \multicolumn{4}{p{0.6\linewidth}}{
        \cellcolor[rgb]{0.95,0.95,0.95}{
            \makecell[{{p{\linewidth}}}]{
                \texttt{\tiny{[P1$\rightarrow$P2]}}
                \texttt{{[}SELF{-}FEEDBACK{]}} \\
\texttt{FORMAT: ok} \\
\texttt{RULES: ok} \\
\texttt{STRATEGY: problem} \\
\texttt{REASON: The move is strategically flawed because it ignores the player's own valuation of canvases (1) in favor of items the other player likely values less, potentially reducing the total combined score compared to claiming the canvases.} \\
            }
        }
    }
     \\ \\

    \theutterance \stepcounter{utterance}  
    & & & \multicolumn{4}{p{0.6\linewidth}}{
        \cellcolor[rgb]{0.9,0.9,0.9}{
            \makecell[{{p{\linewidth}}}]{
                \texttt{\tiny{[GM$\langle$P2]}}
                \texttt{2 canvases} \\
            }
        }
    }
     \\ \\

    \theutterance \stepcounter{utterance}  
    & & & \multicolumn{2}{p{0.3\linewidth}}{
        \cellcolor[rgb]{0.9,0.9,0.9}{
            \makecell[{{p{\linewidth}}}]{
                \texttt{\tiny{[GM$|$GM]}}
                \texttt{continue} \\
            }
        }
    }
    & & \\ \\

    \theutterance \stepcounter{utterance}  
    & & & \multicolumn{2}{p{0.3\linewidth}}{
        \cellcolor[rgb]{0.9,0.9,0.9}{
            \makecell[{{p{\linewidth}}}]{
                \texttt{\tiny{[GM$|$GM]}}
                \texttt{2 canvases} \\
            }
        }
    }
    & & \\ \\

    \theutterance \stepcounter{utterance}  
    & & \multicolumn{4}{p{0.6\linewidth}}{
        \cellcolor[rgb]{0.9,0.9,0.9}{
            \makecell[{{p{\linewidth}}}]{
                \texttt{\tiny{[P1$\langle$GM]}}
                \texttt{2 canvases} \\
            }
        }
    }
    & \\ \\

    \theutterance \stepcounter{utterance}  
    & \multicolumn{4}{p{0.6\linewidth}}{
        \cellcolor[rgb]{0.95,0.95,0.95}{
            \makecell[{{p{\linewidth}}}]{
                \texttt{\tiny{[P2$\rightarrow$P1]}}
                \texttt{{[}DRAFT{]}} \\
\texttt{2 flowers, 3 socks} \\
            }
        }
    }
    & & \\ \\

    \theutterance \stepcounter{utterance}  
    & \multicolumn{4}{p{0.6\linewidth}}{
        \cellcolor[rgb]{0.95,0.95,0.95}{
            \makecell[{{p{\linewidth}}}]{
                \texttt{\tiny{[P2$\rightarrow$P1]}}
                \texttt{{[}SELF{-}FEEDBACK{]}} \\
\texttt{FORMAT: ok} \\
\texttt{RULES: ok} \\
\texttt{STRATEGY: problem} \\
\texttt{REASON: The move is redundant because the player has already proposed "2 flowers, 3 socks" in the previous turn, and repeating it does not advance the negotiation or utilize the remaining message allowance effectively.} \\
            }
        }
    }
    & & \\ \\

    \theutterance \stepcounter{utterance}  
    & \multicolumn{4}{p{0.6\linewidth}}{
        \cellcolor[rgb]{0.9,0.9,0.9}{
            \makecell[{{p{\linewidth}}}]{
                \texttt{\tiny{[P1$\rangle$GM]}}
                \texttt{2 canvases} \\
            }
        }
    }
    & & \\ \\

    \theutterance \stepcounter{utterance}  
    & & & \multicolumn{2}{p{0.3\linewidth}}{
        \cellcolor[rgb]{0.9,0.9,0.9}{
            \makecell[{{p{\linewidth}}}]{
                \texttt{\tiny{[GM$|$GM]}}
                \texttt{continue} \\
            }
        }
    }
    & & \\ \\

    \theutterance \stepcounter{utterance}  
    & & & \multicolumn{2}{p{0.3\linewidth}}{
        \cellcolor[rgb]{0.9,0.9,0.9}{
            \makecell[{{p{\linewidth}}}]{
                \texttt{\tiny{[GM$|$GM]}}
                \texttt{2 canvases} \\
            }
        }
    }
    & & \\ \\

    \theutterance \stepcounter{utterance}  
    & & \multicolumn{4}{p{0.6\linewidth}}{
        \cellcolor[rgb]{0.9,0.9,0.9}{
            \makecell[{{p{\linewidth}}}]{
                \texttt{\tiny{[GM$\rangle$P2]}}
                \texttt{2 canvases} \\
            }
        }
    }
    & \\ \\

    \theutterance \stepcounter{utterance}  
    & & & \multicolumn{4}{p{0.6\linewidth}}{
        \cellcolor[rgb]{0.95,0.95,0.95}{
            \makecell[{{p{\linewidth}}}]{
                \texttt{\tiny{[P1$\rightarrow$P2]}}
                \texttt{{[}DRAFT{]}} \\
\texttt{2 flowers, 3 socks} \\
            }
        }
    }
     \\ \\

    \theutterance \stepcounter{utterance}  
    & & & \multicolumn{4}{p{0.6\linewidth}}{
        \cellcolor[rgb]{0.95,0.95,0.95}{
            \makecell[{{p{\linewidth}}}]{
                \texttt{\tiny{[P1$\rightarrow$P2]}}
                \texttt{{[}SELF{-}FEEDBACK{]}} \\
\texttt{FORMAT: problem} \\
\texttt{RULES: problem} \\
\texttt{STRATEGY: problem} \\
\texttt{REASON: The message contains explanatory text and lacks the required secret proposal format (square brackets), violating both the output format and the rule that non{-}secret messages must not use brackets or include extra commentary.} \\
            }
        }
    }
     \\ \\

    \theutterance \stepcounter{utterance}  
    & & & \multicolumn{4}{p{0.6\linewidth}}{
        \cellcolor[rgb]{0.9,0.9,0.9}{
            \makecell[{{p{\linewidth}}}]{
                \texttt{\tiny{[GM$\langle$P2]}}
                \texttt{problem} \\
\texttt{The draft includes internal monologue and explanatory text rather than a direct message to the other player. It also incorrectly suggests using square brackets in the thought process, which violates the rule that brackets are reserved for the final secret proposal. Furthermore, the strategy is incomplete and does not explicitly aim to maximize the total sum of scores by ensuring complementary proposals.} \\
\\ 
\texttt{RULES} \\
\texttt{1. Do not include any response to the rules or text announcing your message.} \\
\texttt{2. Do not use square brackets when communicating to the other player; use them only for the secret proposal.} \\
\texttt{3. The goal is to maximize the sum of your score and the other player's score.} \\
\texttt{4. Proposals must be complementary (the union of wanted items must not exceed available items).} \\
\\ 
\texttt{STRATEGY} \\
\texttt{To maximize the total score, I should propose a set of items that I value highly while leaving the remaining items available for the other player to claim, assuming they value those remaining items. Since I value socks (3) and flowers (2) more than canvases (2), and you have indicated you value canvases highly, I will claim the flowers and socks. By claiming only these items, I leave the canvases available for you. If you propose to take the canvases, the union is valid (2 flowers + 3 socks + 2 canvases = all items), and we both get positive scores. This approach prioritizes the cooperative goal of utilizing all items efficiently based on our inferred valuations.} \\
\\ 
\texttt{REASON} \\
\texttt{The previous draft failed because it was not a clean message. The refined move is simply stating the desired allocation without brackets or extra text. By proposing "2 flowers, 3 socks", I secure my highest value items (2*1 + 3*2 = 8 points). I leave the 2 canvases open. If the other player proposes to take the canvases, the union is valid, and we both get points. This approach prioritizes the cooperative goal of utilizing all items efficiently based on my known valuations.} \\
\\ 
\texttt{2 flowers, 3 socks} \\
            }
        }
    }
     \\ \\

    \theutterance \stepcounter{utterance}  
    & & & \multicolumn{2}{p{0.3\linewidth}}{
        \cellcolor[rgb]{0.9,0.9,0.9}{
            \makecell[{{p{\linewidth}}}]{
                \texttt{\tiny{[GM$|$GM]}}
                \texttt{continue} \\
            }
        }
    }
    & & \\ \\

    \theutterance \stepcounter{utterance}  
    & & & \multicolumn{2}{p{0.3\linewidth}}{
        \cellcolor[rgb]{0.9,0.9,0.9}{
            \makecell[{{p{\linewidth}}}]{
                \texttt{\tiny{[GM$|$GM]}}
                \texttt{problem} \\
\texttt{The draft includes internal monologue and explanatory text rather than a direct message to the other player. It also incorrectly suggests using square brackets in the thought process, which violates the rule that brackets are reserved for the final secret proposal. Furthermore, the strategy is incomplete and does not explicitly aim to maximize the total sum of scores by ensuring complementary proposals.} \\
\\ 
\texttt{RULES} \\
\texttt{1. Do not include any response to the rules or text announcing your message.} \\
\texttt{2. Do not use square brackets when communicating to the other player; use them only for the secret proposal.} \\
\texttt{3. The goal is to maximize the sum of your score and the other player's score.} \\
\texttt{4. Proposals must be complementary (the union of wanted items must not exceed available items).} \\
\\ 
\texttt{STRATEGY} \\
\texttt{To maximize the total score, I should propose a set of items that I value highly while leaving the remaining items available for the other player to claim, assuming they value those remaining items. Since I value socks (3) and flowers (2) more than canvases (2), and you have indicated you value canvases highly, I will claim the flowers and socks. By claiming only these items, I leave the canvases available for you. If you propose to take the canvases, the union is valid (2 flowers + 3 socks + 2 canvases = all items), and we both get positive scores. This approach prioritizes the cooperative goal of utilizing all items efficiently based on our inferred valuations.} \\
\\ 
\texttt{REASON} \\
\texttt{The previous draft failed because it was not a clean message. The refined move is simply stating the desired allocation without brackets or extra text. By proposing "2 flowers, 3 socks", I secure my highest value items (2*1 + 3*2 = 8 points). I leave the 2 canvases open. If the other player proposes to take the canvases, the union is valid, and we both get points. This approach prioritizes the cooperative goal of utilizing all items efficiently based on my known valuations.} \\
\\ 
\texttt{2 flowers, 3 socks} \\
            }
        }
    }
    & & \\ \\

    \theutterance \stepcounter{utterance}  
    & & \multicolumn{4}{p{0.6\linewidth}}{
        \cellcolor[rgb]{0.9,0.9,0.9}{
            \makecell[{{p{\linewidth}}}]{
                \texttt{\tiny{[P1$\langle$GM]}}
                \texttt{problem} \\
\texttt{The draft includes internal monologue and explanatory text rather than a direct message to the other player. It also incorrectly suggests using square brackets in the thought process, which violates the rule that brackets are reserved for the final secret proposal. Furthermore, the strategy is incomplete and does not explicitly aim to maximize the total sum of scores by ensuring complementary proposals.} \\
\\ 
\texttt{RULES} \\
\texttt{1. Do not include any response to the rules or text announcing your message.} \\
\texttt{2. Do not use square brackets when communicating to the other player; use them only for the secret proposal.} \\
\texttt{3. The goal is to maximize the sum of your score and the other player's score.} \\
\texttt{4. Proposals must be complementary (the union of wanted items must not exceed available items).} \\
\\ 
\texttt{STRATEGY} \\
\texttt{To maximize the total score, I should propose a set of items that I value highly while leaving the remaining items available for the other player to claim, assuming they value those remaining items. Since I value socks (3) and flowers (2) more than canvases (2), and you have indicated you value canvases highly, I will claim the flowers and socks. By claiming only these items, I leave the canvases available for you. If you propose to take the canvases, the union is valid (2 flowers + 3 socks + 2 canvases = all items), and we both get positive scores. This approach prioritizes the cooperative goal of utilizing all items efficiently based on our inferred valuations.} \\
\\ 
\texttt{REASON} \\
\texttt{The previous draft failed because it was not a clean message. The refined move is simply stating the desired allocation without brackets or extra text. By proposing "2 flowers, 3 socks", I secure my highest value items (2*1 + 3*2 = 8 points). I leave the 2 canvases open. If the other player proposes to take the canvases, the union is valid, and we both get points. This approach prioritizes the cooperative goal of utilizing all items efficiently based on my known valuations.} \\
\\ 
\texttt{2 flowers, 3 socks} \\
            }
        }
    }
    & \\ \\

    \theutterance \stepcounter{utterance}  
    & \multicolumn{4}{p{0.6\linewidth}}{
        \cellcolor[rgb]{0.95,0.95,0.95}{
            \makecell[{{p{\linewidth}}}]{
                \texttt{\tiny{[P2$\rightarrow$P1]}}
                \texttt{{[}DRAFT{]}} \\
\texttt{2 flowers, 3 socks} \\
            }
        }
    }
    & & \\ \\

    \theutterance \stepcounter{utterance}  
    & \multicolumn{4}{p{0.6\linewidth}}{
        \cellcolor[rgb]{0.95,0.95,0.95}{
            \makecell[{{p{\linewidth}}}]{
                \texttt{\tiny{[P2$\rightarrow$P1]}}
                \texttt{{[}SELF{-}FEEDBACK{]}} \\
\texttt{FORMAT: ok} \\
\texttt{RULES: ok} \\
\texttt{STRATEGY: problem} \\
\texttt{REASON: The move is redundant as the player has already proposed these exact items in the previous turn, failing to advance the negotiation or utilize the remaining message allowance effectively.} \\
            }
        }
    }
    & & \\ \\

    \theutterance \stepcounter{utterance}  
    & \multicolumn{4}{p{0.6\linewidth}}{
        \cellcolor[rgb]{0.9,0.9,0.9}{
            \makecell[{{p{\linewidth}}}]{
                \texttt{\tiny{[P1$\rangle$GM]}}
                \texttt{2 canvases} \\
            }
        }
    }
    & & \\ \\

    \theutterance \stepcounter{utterance}  
    & & & \multicolumn{2}{p{0.3\linewidth}}{
        \cellcolor[rgb]{0.9,0.9,0.9}{
            \makecell[{{p{\linewidth}}}]{
                \texttt{\tiny{[GM$|$GM]}}
                \texttt{continue} \\
            }
        }
    }
    & & \\ \\

    \theutterance \stepcounter{utterance}  
    & & & \multicolumn{2}{p{0.3\linewidth}}{
        \cellcolor[rgb]{0.9,0.9,0.9}{
            \makecell[{{p{\linewidth}}}]{
                \texttt{\tiny{[GM$|$GM]}}
                \texttt{2 canvases} \\
            }
        }
    }
    & & \\ \\

    \theutterance \stepcounter{utterance}  
    & & \multicolumn{4}{p{0.6\linewidth}}{
        \cellcolor[rgb]{0.9,0.9,0.9}{
            \makecell[{{p{\linewidth}}}]{
                \texttt{\tiny{[GM$\rangle$P2]}}
                \texttt{2 canvases} \\
            }
        }
    }
    & \\ \\

    \theutterance \stepcounter{utterance}  
    & & & \multicolumn{4}{p{0.6\linewidth}}{
        \cellcolor[rgb]{0.95,0.95,0.95}{
            \makecell[{{p{\linewidth}}}]{
                \texttt{\tiny{[P1$\rightarrow$P2]}}
                \texttt{{[}DRAFT{]}} \\
\texttt{2 flowers, 3 socks} \\
            }
        }
    }
     \\ \\

    \theutterance \stepcounter{utterance}  
    & & & \multicolumn{4}{p{0.6\linewidth}}{
        \cellcolor[rgb]{0.95,0.95,0.95}{
            \makecell[{{p{\linewidth}}}]{
                \texttt{\tiny{[P1$\rightarrow$P2]}}
                \texttt{{[}SELF{-}FEEDBACK{]}} \\
\texttt{FORMAT: problem} \\
\texttt{RULES: problem} \\
\texttt{STRATEGY: problem} \\
\texttt{REASON: The answer includes explanatory text and internal monologue instead of a direct message, uses no required tags, and proposes items that are already claimed by the other player, making the proposals non{-}complementary.} \\
            }
        }
    }
     \\ \\

    \theutterance \stepcounter{utterance}  
    & & & \multicolumn{4}{p{0.6\linewidth}}{
        \cellcolor[rgb]{0.9,0.9,0.9}{
            \makecell[{{p{\linewidth}}}]{
                \texttt{\tiny{[GM$\langle$P2]}}
                \texttt{problem} \\
\texttt{The draft includes internal monologue and explanatory text rather than a direct message to the other player. It also incorrectly suggests using square brackets in the thought process, which violates the rule that brackets are reserved for the final secret proposal. Furthermore, the strategy is incomplete and does not explicitly aim to maximize the total sum of scores by ensuring complementary proposals.} \\
\\ 
\texttt{RULES} \\
\texttt{1. Do not include any response to the rules or text announcing your message.} \\
\texttt{2. Do not use square brackets when communicating to the other player; use them only for the secret proposal.} \\
\texttt{3. The goal is to maximize the sum of your score and the other player's score.} \\
\texttt{4. Proposals must be complementary (the union of wanted items must not exceed available items).} \\
\\ 
\texttt{STRATEGY} \\
\texttt{To maximize the total score, I should propose a set of items that I value highly while leaving the remaining items available for the other player to claim, assuming they value those remaining items. Since I value socks (2) and flowers (1) more than canvases (1), and the other player has indicated they value canvases highly, I will claim the flowers and socks. By claiming only these items, I leave the canvases available for the other player. If the other player proposes to take the canvases, the union is valid (2 flowers + 3 socks + 2 canvases = all items), and we both get positive scores. This approach prioritizes the cooperative goal of utilizing all items efficiently based on our inferred valuations.} \\
\\ 
\texttt{REASON} \\
\texttt{The previous draft failed because it was not a clean message. The refined move is simply stating the desired allocation without brackets or extra text. By proposing "2 flowers, 3 socks", I secure my highest value items (2*1 + 3*2 = 8 points). I leave the 2 canvases open. If the other player proposes to take the canvases, the union is valid, and we both get points. This approach prioritizes the cooperative goal of utilizing all items efficiently based on my known valuations.} \\
\\ 
\texttt{2 flowers, 3 socks} \\
            }
        }
    }
     \\ \\

    \theutterance \stepcounter{utterance}  
    & & & \multicolumn{2}{p{0.3\linewidth}}{
        \cellcolor[rgb]{0.9,0.9,0.9}{
            \makecell[{{p{\linewidth}}}]{
                \texttt{\tiny{[GM$|$GM]}}
                \texttt{continue} \\
            }
        }
    }
    & & \\ \\

    \theutterance \stepcounter{utterance}  
    & & & \multicolumn{2}{p{0.3\linewidth}}{
        \cellcolor[rgb]{0.9,0.9,0.9}{
            \makecell[{{p{\linewidth}}}]{
                \texttt{\tiny{[GM$|$GM]}}
                \texttt{problem} \\
\texttt{The draft includes internal monologue and explanatory text rather than a direct message to the other player. It also incorrectly suggests using square brackets in the thought process, which violates the rule that brackets are reserved for the final secret proposal. Furthermore, the strategy is incomplete and does not explicitly aim to maximize the total sum of scores by ensuring complementary proposals.} \\
\\ 
\texttt{RULES} \\
\texttt{1. Do not include any response to the rules or text announcing your message.} \\
\texttt{2. Do not use square brackets when communicating to the other player; use them only for the secret proposal.} \\
\texttt{3. The goal is to maximize the sum of your score and the other player's score.} \\
\texttt{4. Proposals must be complementary (the union of wanted items must not exceed available items).} \\
\\ 
\texttt{STRATEGY} \\
\texttt{To maximize the total score, I should propose a set of items that I value highly while leaving the remaining items available for the other player to claim, assuming they value those remaining items. Since I value socks (2) and flowers (1) more than canvases (1), and the other player has indicated they value canvases highly, I will claim the flowers and socks. By claiming only these items, I leave the canvases available for the other player. If the other player proposes to take the canvases, the union is valid (2 flowers + 3 socks + 2 canvases = all items), and we both get positive scores. This approach prioritizes the cooperative goal of utilizing all items efficiently based on our inferred valuations.} \\
\\ 
\texttt{REASON} \\
\texttt{The previous draft failed because it was not a clean message. The refined move is simply stating the desired allocation without brackets or extra text. By proposing "2 flowers, 3 socks", I secure my highest value items (2*1 + 3*2 = 8 points). I leave the 2 canvases open. If the other player proposes to take the canvases, the union is valid, and we both get points. This approach prioritizes the cooperative goal of utilizing all items efficiently based on my known valuations.} \\
\\ 
\texttt{2 flowers, 3 socks} \\
            }
        }
    }
    & & \\ \\

    \theutterance \stepcounter{utterance}  
    & & \multicolumn{4}{p{0.6\linewidth}}{
        \cellcolor[rgb]{0.9,0.9,0.9}{
            \makecell[{{p{\linewidth}}}]{
                \texttt{\tiny{[P1$\langle$GM]}}
                \texttt{problem} \\
\texttt{The draft includes internal monologue and explanatory text rather than a direct message to the other player. It also incorrectly suggests using square brackets in the thought process, which violates the rule that brackets are reserved for the final secret proposal. Furthermore, the strategy is incomplete and does not explicitly aim to maximize the total sum of scores by ensuring complementary proposals.} \\
\\ 
\texttt{RULES} \\
\texttt{1. Do not include any response to the rules or text announcing your message.} \\
\texttt{2. Do not use square brackets when communicating to the other player; use them only for the secret proposal.} \\
\texttt{3. The goal is to maximize the sum of your score and the other player's score.} \\
\texttt{4. Proposals must be complementary (the union of wanted items must not exceed available items).} \\
\\ 
\texttt{STRATEGY} \\
\texttt{To maximize the total score, I should propose a set of items that I value highly while leaving the remaining items available for the other player to claim, assuming they value those remaining items. Since I value socks (2) and flowers (1) more than canvases (1), and the other player has indicated they value canvases highly, I will claim the flowers and socks. By claiming only these items, I leave the canvases available for the other player. If the other player proposes to take the canvases, the union is valid (2 flowers + 3 socks + 2 canvases = all items), and we both get positive scores. This approach prioritizes the cooperative goal of utilizing all items efficiently based on our inferred valuations.} \\
\\ 
\texttt{REASON} \\
\texttt{The previous draft failed because it was not a clean message. The refined move is simply stating the desired allocation without brackets or extra text. By proposing "2 flowers, 3 socks", I secure my highest value items (2*1 + 3*2 = 8 points). I leave the 2 canvases open. If the other player proposes to take the canvases, the union is valid, and we both get points. This approach prioritizes the cooperative goal of utilizing all items efficiently based on my known valuations.} \\
\\ 
\texttt{2 flowers, 3 socks} \\
\\ 
\\ 
\texttt{The time is up. It is now your turn to submit a secret proposal.} \\
            }
        }
    }
    & \\ \\

    \theutterance \stepcounter{utterance}  
    & \multicolumn{4}{p{0.6\linewidth}}{
        \cellcolor[rgb]{0.95,0.95,0.95}{
            \makecell[{{p{\linewidth}}}]{
                \texttt{\tiny{[P2$\rightarrow$P1]}}
                \texttt{{[}DRAFT{]}} \\
\texttt{{[}Proposal: 2 canvases{]}} \\
            }
        }
    }
    & & \\ \\

    \theutterance \stepcounter{utterance}  
    & \multicolumn{4}{p{0.6\linewidth}}{
        \cellcolor[rgb]{0.95,0.95,0.95}{
            \makecell[{{p{\linewidth}}}]{
                \texttt{\tiny{[P2$\rightarrow$P1]}}
                \texttt{{[}SELF{-}FEEDBACK{]}} \\
\texttt{FORMAT: ok} \\
\texttt{RULES: ok} \\
\texttt{STRATEGY: ok} \\
\texttt{REASON: The proposal correctly claims the high{-}value canvases while leaving the flowers and socks available for the other player, ensuring a complementary and optimal total score.} \\
            }
        }
    }
    & & \\ \\

    \theutterance \stepcounter{utterance}  
    & \multicolumn{4}{p{0.6\linewidth}}{
        \cellcolor[rgb]{0.9,0.9,0.9}{
            \makecell[{{p{\linewidth}}}]{
                \texttt{\tiny{[P1$\rangle$GM]}}
                \texttt{{[}Proposal: 2 canvases{]}} \\
            }
        }
    }
    & & \\ \\

    \theutterance \stepcounter{utterance}  
    & & & \multicolumn{2}{p{0.3\linewidth}}{
        \cellcolor[rgb]{0.9,0.9,0.9}{
            \makecell[{{p{\linewidth}}}]{
                \texttt{\tiny{[GM$|$GM]}}
                \texttt{proposal} \\
            }
        }
    }
    & & \\ \\

    \theutterance \stepcounter{utterance}  
    & & & \multicolumn{2}{p{0.3\linewidth}}{
        \cellcolor[rgb]{0.9,0.9,0.9}{
            \makecell[{{p{\linewidth}}}]{
                \texttt{\tiny{[GM$|$GM]}}
                [0, 0, 2]
            }
        }
    }
    & & \\ \\

    \theutterance \stepcounter{utterance}  
    & & \multicolumn{4}{p{0.6\linewidth}}{
        \cellcolor[rgb]{0.9,0.9,0.9}{
            \makecell[{{p{\linewidth}}}]{
                \texttt{\tiny{[GM$\rangle$P2]}}
                \texttt{The other player has submitted a secret proposal. It is now your turn to submit a proposal of your own.} \\
            }
        }
    }
    & \\ \\

    \theutterance \stepcounter{utterance}  
    & & & \multicolumn{4}{p{0.6\linewidth}}{
        \cellcolor[rgb]{0.95,0.95,0.95}{
            \makecell[{{p{\linewidth}}}]{
                \texttt{\tiny{[P1$\rightarrow$P2]}}
                \texttt{{[}DRAFT{]}} \\
\texttt{{[}Proposal: 2 flowers, 3 socks{]}} \\
            }
        }
    }
     \\ \\

    \theutterance \stepcounter{utterance}  
    & & & \multicolumn{4}{p{0.6\linewidth}}{
        \cellcolor[rgb]{0.95,0.95,0.95}{
            \makecell[{{p{\linewidth}}}]{
                \texttt{\tiny{[P1$\rightarrow$P2]}}
                \texttt{{[}SELF{-}FEEDBACK{]}} \\
\texttt{FORMAT: ok} \\
\texttt{RULES: problem} \\
\texttt{STRATEGY: ok} \\
\texttt{REASON: The answer uses square brackets which triggers a secret proposal instead of a public message, violating the rule that brackets are reserved only for the final secret proposal.} \\
            }
        }
    }
     \\ \\

    \theutterance \stepcounter{utterance}  
    & & & \multicolumn{4}{p{0.6\linewidth}}{
        \cellcolor[rgb]{0.9,0.9,0.9}{
            \makecell[{{p{\linewidth}}}]{
                \texttt{\tiny{[GM$\langle$P2]}}
                \texttt{2 flowers, 3 socks} \\
            }
        }
    }
     \\ \\

    \theutterance \stepcounter{utterance}  
    & & & \multicolumn{2}{p{0.3\linewidth}}{
        \cellcolor[rgb]{0.9,0.9,0.9}{
            \makecell[{{p{\linewidth}}}]{
                \texttt{\tiny{[GM$|$GM]}}
                \texttt{continue} \\
            }
        }
    }
    & & \\ \\

    \theutterance \stepcounter{utterance}  
    & & & \multicolumn{2}{p{0.3\linewidth}}{
        \cellcolor[rgb]{0.9,0.9,0.9}{
            \makecell[{{p{\linewidth}}}]{
                \texttt{\tiny{[GM$|$GM]}}
                \texttt{player was instructed to make a proposal} \\
            }
        }
    }
    & & \\ \\

\end{supertabular}
}

\end{document}
